% Part: sets-functions-relations
% Chapter: relations
% Section: relations-as-sets

\documentclass[../../../include/open-logic-section]{subfiles}

\begin{document}

\olfileid[fa-IR]{sfr}{rel}{set}
\olsection{روابط به‌مثابه مجموعه‌ها}

\begin{explain}
در \olref[sfr][set][imp]{sec}، از چند مجموعهٔ مهم یاد کردیم:
$\Nat$، $\Int$، $\Rat$ و $\Real$. بی‌گمان برخی روابط جالب میان
!!{element}s بعضی از این مجموعه‌ها را به خاطر دارید. برای نمونه، بر
هر یک از این مجموعه‌ها یک \emph{رابطهٔ ترتیب} کاملاً استاندارد وجود
دارد. رابطه‌ای نیز هست با معنای \emph{با یکدیگر یکسان‌بودن}؛ هر شیء
این رابطه را با خودش، و با هیچ چیز دیگری، دارد. روابط جالب دیگری نیز
هستند که با آنها روبه‌رو خواهیم شد، و روابط ممکن از این هم بیشترند.
اما پیش از مرور آنها، نخست نشان می‌دهیم که می‌توان روابط را نوع خاصی
از مجموعه دانست.

برای این منظور، دو مطلب را از \olref[sfr][set][pai]{sec} به یاد آورید.
نخست، مفهوم \emph{زوج مرتب} را یادآوری کنید: با داشتن $a$ و $b$،
می‌توانیم~$\tuple{a, b}$ را تشکیل دهیم. نکتهٔ مهم این است که ترتیب
اعضا در اینجا \emph{واقعاً} اهمیت دارد. پس اگر $a \neq b$، آنگاه
$\tuple{a, b} \neq \tuple{b, a}$. (این را با زوج‌های نامرتب، یعنی
مجموعه‌های $2$-عضوی، مقایسه کنید که در آنها
$\{a, b\}=\{b, a\}$.) دوم، مفهوم \emph{ضرب دکارتی} را به یاد آورید:
اگر $A$ و $B$ مجموعه باشند، می‌توانیم~$A \times B$، یعنی مجموعهٔ همهٔ
زوج‌های $\tuple{x, y}$ را تشکیل دهیم که در آنها $x \in A$ و $y \in B$.
به‌ویژه، $A^{2}= A \times A$ مجموعهٔ همهٔ زوج‌های مرتب از~$A$ است.

اکنون رابطه‌ای خاص روی یک مجموعه را در نظر می‌گیریم: رابطهٔ $<$ روی
مجموعهٔ~$\Nat$ از اعداد طبیعی. مجموعهٔ همهٔ زوج‌های عددی
$\tuple{n, m}$ را در نظر بگیرید که در آنها $n<m$، یعنی
\[
  R=\Setabs{\tuple{n, m}}{n, m \in \Nat \text{ و } n<m}.
\]
میان کوچک‌تر بودن $n$ از $m$ و عضو بودن زوج $\tuple{n, m}$ در $R$
پیوند نزدیکی وجود دارد؛ دقیقاً:
\[
      n<m\text{ اگر و تنها اگر }\tuple{n, m} \in R.
\]
در واقع، بی‌آنکه هیچ اطلاعاتی از دست برود، می‌توانیم مجموعهٔ $R$ را
\emph{همان} رابطهٔ $<$ روی $\Nat$ در نظر بگیریم.

به همین شیوه می‌توانیم برای هر رابطه‌ای میان اعداد، زیرمجموعه‌ای از
$\Nat^{2}$ بسازیم. برعکس، با داشتن هر مجموعهٔ $S \subseteq \Nat^{2}$
از زوج‌های عددی، رابطه‌ای متناظر میان اعداد وجود دارد: رابطه‌ای که
میان $n$ و $m$ برقرار است اگر و تنها اگر $\tuple{n, m} \in S$. این نکته تعریف
زیر را توجیه می‌کند:
\end{explain}

\begin{defn}[رابطهٔ دوتایی]
یک \emph{رابطهٔ دوتایی} روی مجموعهٔ $A$، زیرمجموعه‌ای از~$A^{2}$ است.
اگر $R \subseteq A^{2}$ رابطه‌ای دوتایی روی~$A$ باشد و $x, y \in A$،
گاه $Rxy$ (یا $xRy$) را به جای $\tuple{x, y} \in R$ می‌نویسیم.
\end{defn}

\begin{ex}
  \ollabel{relations}
مجموعهٔ $\Nat^{2}$ از زوج‌های اعداد طبیعی را می‌توان در یک ماتریس
دوبعدی چنین فهرست کرد:
\[
  \begin{array}{ccccc}
  \mathbf{\tuple{ 0,0 }} & \tuple{ 0,1 } &
    \tuple{ 0,2 } & \tuple{ 0,3 } & \ldots\\
  \tuple{ 1,0 } & \mathbf{\tuple{ 1,1 }} &
    \tuple{ 1,2 } & \tuple{ 1,3 } & \ldots\\
  \tuple{ 2,0 } & \tuple{ 2,1 } &
    \mathbf{\tuple{ 2,2 }} & \tuple{ 2,3 } & \ldots\\
  \tuple{ 3,0 } & \tuple{ 3,1 } & \tuple{ 3,2 } &
    \mathbf{\tuple{ 3,3 }} & \ldots\\
  \vdots & \vdots & \vdots & \vdots & \mathbf{\ddots}
  \end{array}
\]
در اینجا درایه‌های روی قطر را پررنگ نوشته‌ایم، زیرا زیرمجموعهٔ
$\Nat^2$ متشکل از زوج‌های روی قطر، یعنی
\[
  \{\tuple{0,0 }, \tuple{ 1,1 }, \tuple{ 2,2 }, \dots\},
  \]
\emph{رابطهٔ همانی روی}~$\Nat$ است. (از آنجا که رابطهٔ همانی
پرکاربرد است، تعریف می‌کنیم $\Id{A}=\Setabs{\tuple{ x,x }}{x \in A}$
برای هر مجموعهٔ $A$.) زیرمجموعهٔ همهٔ زوج‌های
بالای قطر، یعنی
\[
  L = \{\tuple{ 0,1 },\tuple{ 0,2 },\ldots,\tuple{ 1,2 },
  \tuple{ 1,3 }, \dots, \tuple{ 2,3 }, \tuple{ 2,4 },\ldots\},
\]
رابطهٔ \emph{کوچک‌تر از} است؛ یعنی $Lnm$ اگر و تنها اگر $n<m$.
زیرمجموعهٔ زوج‌های پایین قطر، یعنی
\[
  G=\{ \tuple{ 1,0 },\tuple{ 2,0 },\tuple{
    2,1 }, \tuple{ 3,0 },\tuple{ 3,1 },\tuple{ 3,2 }, \dots\},
\]
رابطهٔ \emph{بزرگ‌تر از} است؛ یعنی $Gnm$ اگر و تنها اگر $n>m$.
اجتماع $L$ با $I$ را که می‌توانیم آن را $K=L\cup I$ بنامیم، رابطهٔ
\emph{کوچک‌تر یا مساوی} می‌نامیم: $Knm$ اگر و تنها اگر $n \le m$.
به همین ترتیب، $H=G \cup I$ \emph{رابطهٔ بزرگ‌تر یا مساوی} است. این
روابط $L$، $G$، $K$ و $H$ انواع خاصی از روابط‌اند که \emph{ترتیب}
نامیده می‌شوند. $L$ و $G$ این ویژگی را دارند که هیچ عددی با خودش در
رابطهٔ $L$ یا $G$ نیست (یعنی برای هر $n$، نه $Lnn$ برقرار است و نه
$Gnn$). روابط دارای این ویژگی \emph{غیربازتابی} نامیده می‌شوند و اگر
درعین‌حال ترتیب نیز باشند، \emph{ترتیب اکید} نام دارند.
\end{ex}

\begin{explain}
گرچه ترتیب‌ها و رابطهٔ همانی روابطی مهم و طبیعی‌اند، باید تأکید کرد
که بنا بر تعریف ما، \emph{هر} زیرمجموعهٔ $A^{2}$ رابطه‌ای روی~$A$ است،
هرقدر هم غیرطبیعی یا ساختگی به نظر برسد. به‌ویژه، $\emptyset$ روی هر
مجموعه‌ای یک رابطه است (\emph{رابطهٔ تهی} که هیچ زوجی آن را ندارد)،
و خودِ~$A^{2}$ نیز رابطه‌ای روی~$A$ است (رابطه‌ای که هر زوجی آن را
دارد) و \emph{رابطهٔ جهان‌شمول} نامیده می‌شود. حتی چیزی مانند
$E=\Setabs{\tuple{n, m}}{n>5 \text{ یا } m \times n \ge 34}$ نیز یک
رابطه به شمار می‌آید.
\end{explain}

\begin{prob}
  !!{element}s رابطهٔ $\subseteq$ روی مجموعهٔ
  $\Pow{\{a, b, c\}}$ را فهرست کنید.
\end{prob}

\end{document}
